\part{Independent uses of the mechanism}

\section{What must survive when the equation changes}

The three-dimensional proof does more than produce a terminal bound.  The
datum supplies finite spatial rows of one history; the differentiated
equation generates its temporal and pressure rows; restriction keeps the
overlaps equal; the intersection produces one endpoint; and uniqueness lets
that endpoint restart the evolution.  A genuinely different
three-dimensional producer can preserve those later steps while changing the
estimate that begins them.

Axisymmetric flow without swirl provides the first test.  The physical
vorticity is stretched, but the relative vorticity
\(\omega^\theta/r\) obeys a scalar parabolic law whose maximum principle
controls the stretching coefficient on every finite horizon.  That
datum-side coefficient produces every spatial row, after which the complete
mixed jet, pressure jet, compatible rectangles, endpoint, and restart follow
from the same finite mechanism used in the main proof.  The result is
classical at the level of global smoothness; the work below strengthens it
into a quantitative realization of the complete proof engine.

The later tests change the equation itself.  In viscous--resistive MHD, the
producer must control a coupled velocity and magnetic field while pressure
keeps the difference of their stresses.  Critical hyperdissipation replaces
the adjacent integer Sobolev rows by a fractional Gelfand triple.  The
Navier--Stokes--Voigt equation moves elliptic regularization onto the time
derivative, with the striking consequence that temporal height no longer
increases datum depth.  Each problem is carried from its own datum estimate
through the complete mixed and pressure jets, a common endpoint, and an
explicit same-history restart.  The comparison is made by the proofs, not by
an analogy among their names.

\section{Three-dimensional axisymmetric flow without swirl}
\label{sec:axisymmetric-application}

The next problem restores three-dimensional vortex stretching.  It does not
remove the nonlinearity or pressure, and it permits arbitrarily large smooth
data.  Instead, axial symmetry exposes the precise part of vorticity that the
stretching cannot amplify: the angular vorticity divided by its distance from
the axis.  This scalar becomes the producer for the complete mixed jet.

Write cylindrical coordinates as \((r,\theta,z)\).  A smooth axisymmetric
velocity without swirl has the form
\begin{equation}
 u=u^r(r,z,t)e_r+u^z(r,z,t)e_z,
 \qquad
 \omega:=\nabla\times u=\omega^\theta(r,z,t)e_\theta.
 \label{eq:axisymmetric-form}
\end{equation}
Let \(u_0\in\SSigma(\R^3)\) have this form.  Rotational invariance and local
uniqueness preserve the class on the maximal classical interval.  Smoothness
at the axis gives \(\omega^\theta=O(r)\), so
\begin{equation}
\eta:=\frac{\omega^\theta}{r}
\label{eq:axisymmetric-relative-vorticity}
\end{equation}
extends smoothly across \(r=0\).
Write \(T_{\max}^{\rm axi}\) for the maximal time in this invariant class.

\begin{theorem}[Quantitative axisymmetric no-swirl jet]
\label{thm:quantitative-axisymmetric-jet}
Fix \(p>3\).  Every Schwartz, divergence-free, axisymmetric datum without
swirl generates a unique global smooth no-swirl solution.  For every finite
\(T>0\), every \(n,r,N,M\in\Nzero\), and every
\(q,m\in\mathbb N\) with \(q,m\ge3\), the solution
obeys the datum-side
spatial, Lipschitz, dissipation, velocity-row, pressure-row, rectangle,
endpoint, frequency-tail, and restart estimates
\eqref{eq:axisymmetric-spatial-envelope-bound}--
\eqref{eq:axisymmetric-dissipation-bound},
\eqref{eq:axisymmetric-row-bound},
\eqref{eq:axisymmetric-pressure-bound},
\eqref{eq:axisymmetric-rectangle-bound}, and
\eqref{eq:axisymmetric-endpoint-modulus}--
\eqref{eq:axisymmetric-restart}.
\end{theorem}

\begin{proof}
Directly taking the curl of Navier--Stokes
gives
\begin{align}
 \partial_t\eta+u^r\partial_r\eta+u^z\partial_z\eta
 &=\nu\left(
  \partial_r^2+\frac3r\partial_r+\partial_z^2
 \right)\eta,
 \label{eq:axisymmetric-eta-equation}\\
 \partial_t\omega^\theta+u^r\partial_r\omega^\theta
 +u^z\partial_z\omega^\theta
 -\nu\left(\Delta-\frac1{r^2}\right)\omega^\theta
 &=\frac{u^r}{r}\omega^\theta.
 \label{eq:axisymmetric-vorticity-equation}
\end{align}
These are the classical no-swirl equations.
\bibref{HmidiRousset}\bibref{GallaySverak}

The extra radial term in \eqref{eq:axisymmetric-eta-equation} has the good
axis sign.  Multiplying by \(|\eta|^{a-2}\eta\), integrating with the
three-dimensional measure \(r\,dr\,d\theta\,dz\), and integrating at \(r=0\)
produces the boundary contribution
\begin{equation}
 -\frac{4\pi\nu}{a}\int_{\R}|\eta(0,z,t)|^a\,dz,
 \label{eq:axisymmetric-axis-dissipation}
\end{equation}
in addition to the nonpositive bulk dissipation.  Approximation at
\(a=1\) and the maximum principle then give
\begin{equation}
 \norm{\eta(t)}_{L^a(\R^3)}
 \le\norm{\eta_0}_{L^a(\R^3)}
 \qquad(1\le a\le\infty).
 \label{eq:axisymmetric-eta-producer}
\end{equation}
This scalar producer remains viscous, and the radial diffusion records the
geometry of the axis.

To see what \eqref{eq:axisymmetric-eta-producer} buys, fix \(p>3\).  The
axisymmetric Biot--Savart kernel satisfies
\begin{equation}
 \left|\frac{u^r(x,t)}{r}\right|
 \le C\int_{\R^3}\frac{|\eta(y,t)|}{|x-y|^2}\,dy.
 \label{eq:axisymmetric-biotsavart-pointwise}
\end{equation}
\bibref{HmidiRousset}\bibref{GallaySverak}
Split the integral at a radius \(R>0\).  H\"older's inequality near the
origin and Cauchy--Schwarz outside it give
\begin{equation}
 \norm{u^r(t)/r}_\infty
 \le C_p\left(
 R^{1-3/p}\norm{\eta(t)}_p
 +R^{-1/2}\norm{\eta(t)}_2
 \right).
 \label{eq:axisymmetric-biotsavart-split}
\end{equation}
Optimizing in \(R\) and using
\eqref{eq:axisymmetric-eta-producer} yields
\begin{equation}
 \norm{u^r(t)/r}_\infty
 \le K_p[u_0],
 \qquad
 K_p[u_0]
 :=C_p\norm{\eta_0}_2^{\alpha_p}
       \norm{\eta_0}_p^{1-\alpha_p},
 \label{eq:axisymmetric-stretching-coefficient}
\end{equation}
where
\begin{equation}
 \alpha_p:=\frac{2(p-3)}{3(p-2)}.
 \label{eq:axisymmetric-interpolation-exponent}
\end{equation}
The exponent is forced by balancing the powers
\(R^{1-3/p}\) and \(R^{-1/2}\).  The maximum principle applied to
\eqref{eq:axisymmetric-vorticity-equation} now gives
\begin{equation}
 \norm{\omega(t)}_\infty
 \le W_0e^{K_p[u_0]t},
 \qquad W_0:=\norm{\omega_0}_\infty.
 \label{eq:axisymmetric-vorticity-envelope}
\end{equation}
Thus the genuine three-dimensional stretching coefficient is finite on every
finite horizon and is expressed entirely by the datum.

\subsection{The spatial producer at every order}

Fix an integer \(q\ge3\).  A dyadic decomposition proves the logarithmic
gradient estimate
\begin{equation}
 \norm{\nabla u}_\infty
 \le C_q\left[
  1+\norm{u_0}_2
  +\norm\omega_\infty
   \log\!\left(e+\norm u_{H^q}\right)
 \right].
 \label{eq:axisymmetric-log-gradient}
\end{equation}
Indeed, the negative dyadic shells are bounded by the energy norm
\(\norm{u(t)}_2\le\norm{u_0}_2\).  Each of the first \(J\) nonnegative
shells is bounded by the frequency-localized Biot--Savart multiplier times
\(\norm\omega_\infty\).  The high shells satisfy
\[
 \sum_{j>J}\norm{P_j\nabla u}_\infty
 \le C_q2^{-J(q-5/2)}\norm u_{H^q}.
\]
Choosing \(J\) from \(e+\norm u_{H^q}\) gives
\eqref{eq:axisymmetric-log-gradient}.  This is the usual
Beale--Kato--Majda logarithm, derived here at the exact order used below.
\bibref{BKM}

The complete \(H^q\) energy row gives
\begin{equation}
 \frac d{dt}\norm{u(t)}_{H^q}
 \le C_q\norm{\nabla u(t)}_\infty\norm{u(t)}_{H^q}.
 \label{eq:axisymmetric-high-energy}
\end{equation}
Put
\[
 L_0:=1+\norm{u_0}_2,
 \qquad
 \Psi_K(T):=
 \begin{cases}
  (e^{KT}-1)/K,&K>0,\\
  T,&K=0,
 \end{cases}
\]
\begin{equation}
 A_p(T):=W_0\Psi_{K_p[u_0]}(T),
 \qquad
 Y_{q,0}:=\log\!\left(e+\norm{u_0}_{H^q}\right).
 \label{eq:axisymmetric-envelope-data}
\end{equation}
For
\(Y_q(t):=\log(e+\norm{u(t)}_{H^q})\), equations
\eqref{eq:axisymmetric-vorticity-envelope}--
\eqref{eq:axisymmetric-high-energy} give
\[
 Y_q'(t)
 \le C_q\left(L_0+W_0e^{K_p[u_0]t}Y_q(t)\right).
\]
Gronwall's inequality yields the completely datum-side quantities
\begin{equation}
 \overline Y_q(T)
 :=\bigl(Y_{q,0}+C_qL_0T\bigr)e^{C_qA_p(T)},
 \qquad
 \mathsf S_q^{\rm axi}(T):=e^{\overline Y_q(T)},
 \label{eq:axisymmetric-spatial-envelope}
\end{equation}
and the bounds
\begin{align}
 \sup_{0\le t\le T}\norm{u(t)}_{H^q}
 &\le\mathsf S_q^{\rm axi}(T),
 \label{eq:axisymmetric-spatial-envelope-bound}\\
 \int_0^T\norm{\nabla u(t)}_\infty\,dt
 &\le\mathsf J_q^{\rm axi}(T)
 :=C_q\left[L_0T+A_p(T)\overline Y_q(T)\right].
 \label{eq:axisymmetric-lipschitz-action}
\end{align}
No norm of the evolved solution remains on either right-hand side.

The dissipative companion is also explicit.  The squared high-order energy
row and \eqref{eq:axisymmetric-lipschitz-action} give
\begin{equation}
 \int_0^T\norm{u(t)}_{H^{q+1}}^2\,dt
 \le\mathsf D_q^{\rm axi}(T)^2,
 \label{eq:axisymmetric-dissipation-bound}
\end{equation}
where, after enlarging the norm-equivalence constant \(C_q\), one may take
\begin{equation}
 \mathsf D_q^{\rm axi}(T)^2
 :=C_q\left[
 T\norm{u_0}_2^2
 +\frac{\norm{u_0}_{H^q}^2
        +2C_q\mathsf S_q^{\rm axi}(T)^2
          \mathsf J_q^{\rm axi}(T)}{\nu}
 \right].
 \label{eq:axisymmetric-dissipation-envelope}
\end{equation}

\subsection{The axisymmetric mixed and pressure jets}

The producer has now paid all spatial orders.  The rest of the construction
returns to the original three-dimensional equation.  Put
\(V_n:=\partial_t^nu\), set \(\sigma_r:=\max\{2,r+1\}\), and define
\begin{equation}
 \mathsf U_{0,r}^{\rm axi}(T)
 :=\mathsf S_{\max\{3,r\}}^{\rm axi}(T).
 \label{eq:axisymmetric-row-base}
\end{equation}
With the product constants from \eqref{eq:row-product-constant}, define
\begin{equation}
 \mathsf U_{n+1,r}^{\rm axi}(T)
 :=\nu\mathsf U_{n,r+2}^{\rm axi}(T)
 +A_r\sum_{a=0}^{n}\binom na
 \mathsf U_{a,\sigma_r}^{\rm axi}(T)
 \mathsf U_{n-a,\sigma_r}^{\rm axi}(T).
 \label{eq:axisymmetric-row-recursion}
\end{equation}
The projected momentum recurrence gives
\begin{equation}
 \sup_{0\le t\le T}\norm{V_n(t)}_{H^r}
 \le\mathsf U_{n,r}^{\rm axi}(T).
 \label{eq:axisymmetric-row-bound}
\end{equation}
For the natural negative row, put
\begin{equation}
 \mathsf U_{n+1,-1}^{\rm axi}(T)
 :=\nu\mathsf U_{n,1}^{\rm axi}(T)
 +A_{-1}\sum_{a=0}^{n}\binom na
 \mathsf U_{a,1}^{\rm axi}(T)
 \mathsf U_{n-a,1}^{\rm axi}(T).
 \label{eq:axisymmetric-negative-row}
\end{equation}

The normalized pressure rows are
\begin{equation}
 Q_n:=\partial_t^np
 =R_iR_j\sum_{a=0}^{n}\binom na
  (V_a)_i(V_{n-a})_j.
 \label{eq:axisymmetric-pressure-formula}
\end{equation}
Accordingly, they obey
\begin{equation}
 \mathsf Q_{n,r}^{\rm axi}(T)
 :=A_r\sum_{a=0}^{n}\binom na
 \mathsf U_{a,\sigma_r}^{\rm axi}(T)
 \mathsf U_{n-a,\sigma_r}^{\rm axi}(T),
 \label{eq:axisymmetric-pressure-envelope}
\end{equation}
and hence
\begin{equation}
 \sup_{0\le t\le T}\norm{\partial_t^np(t)}_{H^r}
 \le\mathsf Q_{n,r}^{\rm axi}(T).
 \label{eq:axisymmetric-pressure-bound}
\end{equation}
For every \(N,M\in\Nzero\), define
\begin{equation}
 \mathfrak R_{N,M}^{\rm axi}(T)
 :=T\sum_{n=0}^{N}\left[
  \bigl(\mathsf U_{n,M+1}^{\rm axi}(T)\bigr)^2
  +\bigl(\mathsf U_{n+1,M-1}^{\rm axi}(T)\bigr)^2
 \right],
 \label{eq:axisymmetric-rectangle-envelope}
\end{equation}
using \eqref{eq:axisymmetric-negative-row} when \(M=0\).  Measure the actual
axisymmetric history by
\[
 \mathcal N_{N,M}^{\rm axi}(u;T)
 :=\sum_{n=0}^{N}\norm{V_n}_{L^2(0,T;H^{M+1})}^2
   +\sum_{n=0}^{N}\norm{V_{n+1}}_{L^2(0,T;H^{M-1})}^2.
\]
It satisfies
\begin{equation}
 \mathcal N_{N,M}^{\rm axi}(u;T)
 \le\mathfrak R_{N,M}^{\rm axi}(T).
 \label{eq:axisymmetric-rectangle-bound}
\end{equation}
Expanding this finite recursion uses datum derivatives through order at most
\(\max\{3,M+2N+1\}\), together with \(\eta_0\) in the two displayed
Lebesgue spaces.

Every rectangle is a restriction of the same axisymmetric velocity and
normalized pressure.  At a chosen horizon, the endpoint moduli are
\begin{align}
 \norm{V_n(t)-V_n(T)}_{H^m}
 &\le(T-t)\mathsf U_{n+1,m}^{\rm axi}(T),
 \label{eq:axisymmetric-endpoint-modulus}\\
 \norm{Q_n(t)-Q_n(T)}_{H^r}
 &\le(T-t)\mathsf Q_{n+1,r}^{\rm axi}(T).
 \label{eq:axisymmetric-pressure-modulus}
\end{align}
The critical frequency tail has the explicit form
\begin{equation}
 \sup_{0\le t\le T}H_{>K}(t)
 \le\frac12K^{1-2m}\mathsf S_m^{\rm axi}(T)^2,
 \qquad m\ge3, K\ge1.
 \label{eq:axisymmetric-tail}
\end{equation}

For every \(0<T<T_{\max}^{\rm axi}\), the local \(H^3\) construction and
\eqref{eq:axisymmetric-spatial-envelope-bound} give a same-history restart of
length
\begin{equation}
 \delta_T^{\rm axi}
 \ge\min\left\{1,
 c_\nu\bigl(1+\mathsf S_3^{\rm axi}(T)\bigr)^{-2}
 \right\}.
 \label{eq:axisymmetric-restart}
\end{equation}

If \(T_{\max}^{\rm axi}\) were finite, apply
the estimates first on \([0,\tau]\), \(\tau<T_{\max}^{\rm axi}\), with the
right sides evaluated at \(T_{\max}^{\rm axi}\), and then let
\(\tau\uparrow T_{\max}^{\rm axi}\).  The compatible
rectangles give one \(H^\infty\) endpoint and the complete pressure jet.  At
that endpoint the limiting rows obey
\[
 V_{n+1,*}
 =\nu\Delta V_{n,*}
  -\sum_{a=0}^{n}\binom na
   \PP(V_{a,*}\cdot\nabla V_{n-a,*}),
\]
together with the normalized limit of
\eqref{eq:axisymmetric-pressure-formula}.  The local theorem restarts it for
the positive time supplied by \eqref{eq:axisymmetric-restart} with
\(T=T_{\max}^{\rm axi}\).
The adjacent-row trace proof, the normalized pressure recurrence, and the
local \(H^3\) construction now apply directly to these axisymmetric
rectangles.  The right local jet agrees inductively with the left endpoint
jet, the joined curve is smooth through \(T_{\max}^{\rm axi}\), and uniqueness thereafter
preserves axial symmetry and zero swirl.  Finite maximality is impossible.
This recovers global smoothness for every Schwartz axisymmetric no-swirl datum
in the stated class, now with an explicit mixed jet, pressure jet, frequency
tail, endpoint modulus, and restart time.  The global regularity of this class
is classical; the result here is its complete realization through the paper's
quantitative producer--intersection mechanism.
\end{proof}

\section{A coupled three-dimensional fluid and magnetic field}
\label{sec:mhd-application}

The previous application still evolved one vector field.  A stronger test is
to let a second divergence-free field enter the transport itself.  In
Alfv\'enic units, viscous--resistive incompressible magnetohydrodynamics on
\(\R^3\) reads
\begin{align}
 \partial_tu-\nu\Delta u+(u\cdot\nabla)u-(b\cdot\nabla)b
 +\nabla\Pi&=0,
 \label{eq:mhd-momentum}\\
 \partial_tb-\eta\Delta b+(u\cdot\nabla)b-(b\cdot\nabla)u&=0,
 \label{eq:mhd-induction}\\
 \nabla\cdot u=\nabla\cdot b&=0.
 \label{eq:mhd-divergence}
\end{align}
Here \(\nu,\eta>0\), and
\(\Pi=p+|b|^2/2\) is total pressure.  The magnetic field can transfer
energy to the velocity and take it back.  The producer must control that
exchange as one coupled state; the endpoint must carry two compatible jets;
and the pressure must retain the difference between their quadratic
stresses.

Let \(\dot\Delta_j\) be homogeneous Littlewood--Paley blocks and write
\begin{equation}
 \norm f_{\dot B^s_{2,1}}
 :=\sum_{j\in\mathbb Z}2^{js}\norm{\dot\Delta_jf}_2.
 \label{eq:mhd-besov-norm}
\end{equation}
For a time norm with a tilde, the time norm is taken inside this dyadic sum.
Schwartz data satisfy the usual homogeneous low-frequency normalization.
The critical small-data theory in these spaces is classical.
\bibref{MiaoYuan}

\begin{theorem}[Quantitative coupled MHD jet]
\label{thm:quantitative-mhd-jet}
Let \(u_0,b_0\in\mathcal S(\R^3)\) be divergence free and put
\begin{equation}
 \mu:=\min\{\nu,\eta\},
 \qquad
 A_0:=\norm{u_0}_{\dot B^{1/2}_{2,1}}
      +\norm{b_0}_{\dot B^{1/2}_{2,1}}.
 \label{eq:mhd-data-size}
\end{equation}
There are fixed dyadic constants \(c_h,C_B>0\) such that, if
\begin{equation}
 A_0\le\varepsilon_{\rm MHD}\mu,
 \qquad
 \varepsilon_{\rm MHD}:=\frac{c_h}{4C_B},
 \label{eq:mhd-smallness}
\end{equation}
then \eqref{eq:mhd-momentum}--\eqref{eq:mhd-divergence} has a unique global
smooth solution.  Every finite spatial and temporal order obeys the explicit
coupled producer, velocity, magnetic, total-pressure, mechanical-pressure,
rectangle, endpoint, tail, datum-depth, and restart estimates
\eqref{eq:mhd-critical-producer}--\eqref{eq:mhd-restart} below.
\end{theorem}

\begin{proof}
The first cancellation already shows why \((u,b)\) must be kept together.
Pair \eqref{eq:mhd-momentum} with \(u\), pair
\eqref{eq:mhd-induction} with \(b\), and add.  Incompressibility removes
the two terms transported by \(u\), while integration by parts gives
\[
 -\pair{b\cdot\nabla b}{u}
 -\pair{b\cdot\nabla u}{b}=0.
\]
Thus
\begin{equation}
 \frac12\frac d{dt}
 \left(\norm u_2^2+\norm b_2^2\right)
 +\nu\norm{\nabla u}_2^2+\eta\norm{\nabla b}_2^2=0.
 \label{eq:mhd-energy}
\end{equation}

The energy identity alone does not control the Lipschitz action needed by
the high Sobolev rows.  The critical Besov producer does.  For a finite
horizon define
\begin{align}
 \mathcal A(T)
 &:=\norm u_{\widetilde L^\infty(0,T;\dot B^{1/2}_{2,1})}
   +\norm b_{\widetilde L^\infty(0,T;\dot B^{1/2}_{2,1})},
 \label{eq:mhd-critical-size}\\
 \mathcal D(T)
 &:=\nu\norm u_{L^1(0,T;\dot B^{5/2}_{2,1})}
   +\eta\norm b_{L^1(0,T;\dot B^{5/2}_{2,1})}.
 \label{eq:mhd-critical-dissipation}
\end{align}
The product estimate that bears the weight is
\begin{equation}
 \norm{f\cdot\nabla g}_{\dot B^{1/2}_{2,1}}
 \le C_B\norm f_{\dot B^{1/2}_{2,1}}
          \norm g_{\dot B^{5/2}_{2,1}}.
 \label{eq:mhd-critical-product}
\end{equation}
To see the derivative count, decompose
\[
 f\cdot\nabla g
 =T_f(\nabla g)+T_{\nabla g}f+R(f,\nabla g).
\]
The three paraproduct bounds, together with
\(\dot B^{3/2}_{2,1}\hookrightarrow L^\infty\), give
\eqref{eq:mhd-critical-product}; the endpoint summation is paid by the
Littlewood--Paley \(\ell^1\) index.
\bibref{BCD}

Apply the dyadic heat estimate to the projected momentum equation and the
induction equation.  Summing over \(j\) and using
\eqref{eq:mhd-critical-product} gives
\begin{equation}
 \mathcal A(T)+c_h\mathcal D(T)
 \le A_0+\frac{C_B}{\mu}\mathcal A(T)\mathcal D(T).
 \label{eq:mhd-bootstrap-inequality}
\end{equation}
On an interval where \(\mathcal A(T)\le2A_0\), the smallness assumption
\eqref{eq:mhd-smallness} absorbs the last term and improves the bound.
Continuity closes the bootstrap:
\begin{equation}
 \mathcal A(T)+\frac{c_h}{2}\mathcal D(T)
 \le A_0
 \qquad(T<T_{\max}).
 \label{eq:mhd-critical-producer}
\end{equation}
If \(C_E\) is the embedding constant in
\(\norm{\nabla f}_\infty\le C_E\norm f_{\dot B^{5/2}_{2,1}}\), then
\begin{equation}
 \int_0^T\left(\norm{\nabla u(t)}_\infty
                  +\norm{\nabla b(t)}_\infty\right)dt
 \le J_0:=\frac{2C_EA_0}{c_h\mu}.
 \label{eq:mhd-lipschitz-action}
\end{equation}
The right side no longer depends on the horizon.

Fix an integer \(q\ge3\) and set
\begin{equation}
 Z_q(t):=\norm{u(t)}_{H^q}^2+\norm{b(t)}_{H^q}^2,
 \qquad
 L(t):=\norm{\nabla u(t)}_\infty+\norm{\nabla b(t)}_\infty.
 \label{eq:mhd-high-variables}
\end{equation}
After \(\partial^\gamma\), \(|\gamma|\le q\), the leading
\(u\)-transport terms vanish.  The two leading \(b\)-terms cancel against
each other just as they did in \eqref{eq:mhd-energy}.  Every remaining term
is a commutator, and
\begin{equation}
 \norm{[\partial^\gamma,f\cdot\nabla]g}_2
 \le C_q\left(
  \norm{\nabla f}_\infty\norm g_{H^q}
  +\norm f_{H^q}\norm{\nabla g}_\infty
 \right).
 \label{eq:mhd-commutator}
\end{equation}
Summing the two equations gives
\begin{equation}
 \frac12Z_q'(t)
 +\nu\norm{\nabla u}_{H^q}^2
 +\eta\norm{\nabla b}_{H^q}^2
 \le C_qL(t)Z_q(t).
 \label{eq:mhd-high-energy}
\end{equation}
Define
\begin{equation}
 \mathsf S_q^{\rm MHD}
 :=Z_q(0)^{1/2}e^{C_qJ_0}.
 \label{eq:mhd-spatial-envelope}
\end{equation}
After enlarging \(C_q\) to absorb the harmless factor from
\eqref{eq:mhd-high-energy}, one obtains
\begin{equation}
 \sup_{0\le t\le T}Z_q(t)^{1/2}
 \le\mathsf S_q^{\rm MHD}.
 \label{eq:mhd-spatial-bound}
\end{equation}
The integrated companion is
\begin{equation}
 \bigl(\mathsf D_q^{\rm MHD}(T)\bigr)^2
 :=T\bigl(\mathsf S_q^{\rm MHD}\bigr)^2
 +\frac{\frac12Z_q(0)+C_qJ_0(\mathsf S_q^{\rm MHD})^2}{\mu},
 \label{eq:mhd-dissipation-envelope}
\end{equation}
for which
\begin{equation}
 \norm u_{L^2(0,T;H^{q+1})}^2
 +\norm b_{L^2(0,T;H^{q+1})}^2
 \le\bigl(\mathsf D_q^{\rm MHD}(T)\bigr)^2.
 \label{eq:mhd-dissipation-bound}
\end{equation}
The bound \eqref{eq:mhd-spatial-bound}, first at one \(q>5/2\) and then at
every finite \(q\), prevents finite-time loss of the classical coupled
solution.

\subsection{The coupled mixed jet}

Write
\begin{equation}
 V_n:=\partial_t^nu,
 \qquad W_n:=\partial_t^nb.
 \label{eq:mhd-row-notation}
\end{equation}
Differentiating the projected equations gives
\begin{align}
 V_{n+1}
 &={}\nu\Delta V_n
 -\sum_{a=0}^n\binom na\PP\left[
   (V_a\cdot\nabla)V_{n-a}
   -(W_a\cdot\nabla)W_{n-a}
 \right],
 \label{eq:mhd-velocity-recurrence}\\
 W_{n+1}
 &={}\eta\Delta W_n
 -\sum_{a=0}^n\binom na\left[
   (V_a\cdot\nabla)W_{n-a}
   -(W_a\cdot\nabla)V_{n-a}
 \right].
 \label{eq:mhd-magnetic-recurrence}
\end{align}
For \(r\ge0\), let \(\sigma_r:=\max\{2,r+1\}\) and fix \(A_r\) so that
\begin{equation}
 \norm{\PP(f\cdot\nabla g)}_{H^r}
 +\norm{f\cdot\nabla g}_{H^r}
 \le A_r\norm f_{H^{\sigma_r}}\norm g_{H^{\sigma_r}}.
 \label{eq:mhd-row-product}
\end{equation}
Set
\begin{equation}
 \mathsf U_{0,r}^{\rm MHD}
 =\mathsf B_{0,r}^{\rm MHD}
 :=\mathsf S_{\max\{3,r\}}^{\rm MHD}.
 \label{eq:mhd-row-base}
\end{equation}
Define the two finite recursions
\begin{align}
 \mathsf U_{n+1,r}^{\rm MHD}
 &:={}\nu\mathsf U_{n,r+2}^{\rm MHD}
 +A_r\sum_{a=0}^n\binom na\left(
  \mathsf U_{a,\sigma_r}^{\rm MHD}
  \mathsf U_{n-a,\sigma_r}^{\rm MHD}
  +\mathsf B_{a,\sigma_r}^{\rm MHD}
  \mathsf B_{n-a,\sigma_r}^{\rm MHD}
 \right),
 \label{eq:mhd-velocity-envelope}\\
 \mathsf B_{n+1,r}^{\rm MHD}
 &:={}\eta\mathsf B_{n,r+2}^{\rm MHD}
 +A_r\sum_{a=0}^n\binom na\left(
  \mathsf U_{a,\sigma_r}^{\rm MHD}
  \mathsf B_{n-a,\sigma_r}^{\rm MHD}
  +\mathsf B_{a,\sigma_r}^{\rm MHD}
  \mathsf U_{n-a,\sigma_r}^{\rm MHD}
 \right).
 \label{eq:mhd-magnetic-envelope}
\end{align}
They give
\begin{align}
 \sup_{0\le t\le T}\norm{V_n(t)}_{H^r}
 &\le\mathsf U_{n,r}^{\rm MHD},
 \label{eq:mhd-velocity-bound}\\
 \sup_{0\le t\le T}\norm{W_n(t)}_{H^r}
 &\le\mathsf B_{n,r}^{\rm MHD}.
 \label{eq:mhd-magnetic-bound}
\end{align}
At the natural negative row, fix \(A_{-1}\) so that
\begin{equation}
 \norm{\PP(f\cdot\nabla g)}_{H^{-1}}
 +\norm{f\cdot\nabla g}_{H^{-1}}
 \le A_{-1}\norm f_{H^1}\norm g_{H^1},
 \label{eq:mhd-negative-product}
\end{equation}
for the divergence-free factors used here, because
\(f\cdot\nabla g=\nabla\cdot(f\otimes g)\),
\(\norm{f\otimes g}_2\le\norm f_4\norm g_4\), and \(\PP\) is bounded on
\(H^{-1}\).  Put
\begin{equation}
 \mathsf U_{0,-1}^{\rm MHD}
 =\mathsf B_{0,-1}^{\rm MHD}:=\mathsf S_3^{\rm MHD}
 \label{eq:mhd-negative-base}
\end{equation}
and define
\begin{align}
 \mathsf U_{n+1,-1}^{\rm MHD}
 &:={}\nu\mathsf U_{n,1}^{\rm MHD}
 +A_{-1}\sum_{a=0}^n\binom na\left(
  \mathsf U_{a,1}^{\rm MHD}\mathsf U_{n-a,1}^{\rm MHD}
  +\mathsf B_{a,1}^{\rm MHD}\mathsf B_{n-a,1}^{\rm MHD}
 \right),
 \label{eq:mhd-negative-velocity-envelope}\\
 \mathsf B_{n+1,-1}^{\rm MHD}
 &:={}\eta\mathsf B_{n,1}^{\rm MHD}
 +A_{-1}\sum_{a=0}^n\binom na\left(
  \mathsf U_{a,1}^{\rm MHD}\mathsf B_{n-a,1}^{\rm MHD}
  +\mathsf B_{a,1}^{\rm MHD}\mathsf U_{n-a,1}^{\rm MHD}
 \right).
 \label{eq:mhd-negative-magnetic-envelope}
\end{align}
Equations \eqref{eq:mhd-velocity-recurrence}--
\eqref{eq:mhd-negative-product} give
\begin{equation}
 \sup_{0\le t\le T}
 \left(\norm{V_n(t)}_{H^{-1}}+\norm{W_n(t)}_{H^{-1}}\right)
 \le\mathsf U_{n,-1}^{\rm MHD}+\mathsf B_{n,-1}^{\rm MHD}.
 \label{eq:mhd-negative-bound}
\end{equation}
This supplies the lower side of the \(M=0\) rectangle with the projected
transport term retained.

\subsection{Both pressures and every finite rectangle}

The normalized total pressure is
\begin{equation}
 \Pi=R_iR_j(u_i u_j-b_i b_j).
 \label{eq:mhd-total-pressure}
\end{equation}
For \(Q_n:=\partial_t^n\Pi\),
\begin{equation}
 Q_n=R_iR_j\sum_{a=0}^n\binom na\left[
  (V_a)_i(V_{n-a})_j-(W_a)_i(W_{n-a})_j
 \right].
 \label{eq:mhd-total-pressure-recurrence}
\end{equation}
Let
\begin{equation}
 \tau_0:=1,\qquad \tau_1:=2,\qquad \tau_r:=r\quad(r\ge2),
 \label{eq:mhd-pressure-indices}
\end{equation}
and fix \(P_r\) so that
\begin{equation}
 \norm{R_iR_j(f_i g_j)}_{H^r}+\norm{f_i g_j}_{H^r}
 \le P_r\norm f_{H^{\tau_r}}\norm g_{H^{\tau_r}}.
 \label{eq:mhd-pressure-product}
\end{equation}
Then
\begin{equation}
 \mathsf Q_{n,r}^{\rm MHD}
 :=P_r\sum_{a=0}^n\binom na\left(
  \mathsf U_{a,\tau_r}^{\rm MHD}
  \mathsf U_{n-a,\tau_r}^{\rm MHD}
  +\mathsf B_{a,\tau_r}^{\rm MHD}
  \mathsf B_{n-a,\tau_r}^{\rm MHD}
 \right)
 \label{eq:mhd-total-pressure-envelope}
\end{equation}
obeys
\begin{equation}
 \sup_{0\le t\le T}\norm{Q_n(t)}_{H^r}
 \le\mathsf Q_{n,r}^{\rm MHD}.
 \label{eq:mhd-total-pressure-bound}
\end{equation}
The mechanical pressure is not discarded.  Since
\(p=\Pi-|b|^2/2\), its row is
\begin{equation}
 \partial_t^np
 =Q_n-\frac12\sum_{a=0}^n\binom na W_a\cdot W_{n-a},
 \label{eq:mhd-mechanical-pressure-recurrence}
\end{equation}
Define its explicit companion envelope by
\begin{equation}
 \mathsf P_{n,r}^{\rm mech}(T)
 :=\mathsf Q_{n,r}^{\rm MHD}
 +\frac{P_r}{2}\sum_{a=0}^n\binom na
  \mathsf B_{a,\tau_r}^{\rm MHD}
  \mathsf B_{n-a,\tau_r}^{\rm MHD}.
 \label{eq:mhd-mechanical-pressure-envelope}
\end{equation}
Then
\begin{equation}
 \sup_{0\le t\le T}\norm{\partial_t^np(t)}_{H^r}
 \le\mathsf P_{n,r}^{\rm mech}(T).
 \label{eq:mhd-mechanical-pressure-bound}
\end{equation}

For \(N,M\in\Nzero\), define
\begin{align}
 \mathcal N_{N,M}^{\rm MHD}(u,b;T)
 :=\sum_{n=0}^N\bigl[
 &\norm{V_n}_{L^2(0,T;H^{M+1})}^2
 +\norm{W_n}_{L^2(0,T;H^{M+1})}^2
 \notag\\
 &+\norm{V_{n+1}}_{L^2(0,T;H^{M-1})}^2
 +\norm{W_{n+1}}_{L^2(0,T;H^{M-1})}^2
 \bigr].
 \label{eq:mhd-rectangle-norm}
\end{align}
The uniform rows give the datum-side enclosure
\begin{align}
 \mathcal N_{N,M}^{\rm MHD}(u,b;T)
 &\le\mathfrak R_{N,M}^{\rm MHD}(T),
 \label{eq:mhd-rectangle-bound}\\
 \mathfrak R_{N,M}^{\rm MHD}(T)
 :=T\sum_{n=0}^N\bigl[
 &(\mathsf U_{n,M+1}^{\rm MHD})^2
 +(\mathsf B_{n,M+1}^{\rm MHD})^2
 \notag\\
 &+(\mathsf U_{n+1,M-1}^{\rm MHD})^2
 +(\mathsf B_{n+1,M-1}^{\rm MHD})^2
 \bigr].
 \label{eq:mhd-rectangle-envelope}
\end{align}
The negative rows just defined are used when \(M=0\).  The pressure companion
is
\begin{equation}
 \sum_{n=0}^N\left[
  \norm{Q_n}_{L^2(0,T;H^M)}^2
  +\norm{\partial_t^np}_{L^2(0,T;H^M)}^2
 \right]
 \le T\sum_{n=0}^N\left[
  (\mathsf Q_{n,M}^{\rm MHD})^2
  +(\mathsf P_{n,M}^{\rm mech}(T))^2
 \right].
 \label{eq:mhd-pressure-rectangle}
\end{equation}
Every rectangle is a restriction of the same coupled history, so all
overlaps agree without choosing either field again.

The velocity and magnetic row \((n,r)\) reads the data through no order above
\begin{equation}
 d_{n,r}^{u,b}:=\max\{3,r+2n\}.
 \label{eq:mhd-field-ledger}
\end{equation}
The total or mechanical pressure row reads no higher than
\begin{equation}
 d_{n,r}^{p}:=\max\{3,\tau_r+2n\},
 \label{eq:mhd-pressure-ledger}
\end{equation}
and the rectangle \((N,M)\) reads no higher than
\begin{equation}
 d_{N,M}^{\rm rect}:=\max\{3,M+2N+1\}.
 \label{eq:mhd-rectangle-ledger}
\end{equation}

\subsection{Coupled endpoint and restart}

For every \(n,m\in\Nzero\), Banach-valued integration gives
\begin{align}
 \norm{V_n(t)-V_n(T)}_{H^m}
 &\le(T-t)\mathsf U_{n+1,m}^{\rm MHD},
 \label{eq:mhd-velocity-endpoint}\\
 \norm{W_n(t)-W_n(T)}_{H^m}
 &\le(T-t)\mathsf B_{n+1,m}^{\rm MHD},
 \label{eq:mhd-magnetic-endpoint}\\
 \norm{Q_n(t)-Q_n(T)}_{H^m}
 &\le(T-t)\mathsf Q_{n+1,m}^{\rm MHD},
 \label{eq:mhd-pressure-endpoint}\\
 \norm{\partial_t^np(t)-\partial_t^np(T)}_{H^m}
 &\le(T-t)\mathsf P_{n+1,m}^{\rm mech}(T).
 \label{eq:mhd-mechanical-pressure-endpoint}
\end{align}
The two field moduli read the data through no order above
\begin{equation}
 \max\{3,m+2n+2\}.
 \label{eq:mhd-field-endpoint-ledger}
\end{equation}
The datum depth of both pressure moduli is at most
\begin{equation}
 \max\{3,\tau_m+2n+2\}.
 \label{eq:mhd-pressure-endpoint-ledger}
\end{equation}

For every temporal row define the coupled critical height and its
high-frequency tail by
\begin{align}
 \mathcal H_n^{\rm MHD}(t)
 :=\frac12\int_{\R^3}|\xi|
 \left(|\widehat V_n(\xi,t)|^2+|\widehat W_n(\xi,t)|^2\right)d\xi,
 \label{eq:mhd-critical-height}\\
 \mathcal H_{n,>K}^{\rm MHD}(t)
 :=\frac12\int_{|\xi|>K}|\xi|
 \left(|\widehat V_n|^2+|\widehat W_n|^2\right)d\xi.
 \label{eq:mhd-critical-tail}
\end{align}
Thus \(\mathcal H_0^{\rm MHD}\) is the base coupled height.  The row
envelopes give
\begin{align}
 \sup_{0\le t\le T}\mathcal H_n^{\rm MHD}(t)
 &\le\frac12\left[
  (\mathsf U_{n,1}^{\rm MHD})^2
  +(\mathsf B_{n,1}^{\rm MHD})^2
 \right],
 \label{eq:mhd-critical-height-bound}\\
 \sup_{0\le t\le T}\mathcal H_{n,>K}^{\rm MHD}(t)
 &\le\frac12K^{1-2m}\left[
  (\mathsf U_{n,m}^{\rm MHD})^2
  +(\mathsf B_{n,m}^{\rm MHD})^2
 \right],
 \qquad m\ge3,\ K\ge1.
 \label{eq:mhd-frequency-tail}
\end{align}

It remains to cross a hypothetical finite terminal time with the coupled
state.  Use the linear constant \(C_{\rm lin}\) from
\eqref{eq:explicit-linear-local-constant} and define
\begin{equation}
 C_{\rm sys}(\nu,\eta)
 :=\max\{C_{\rm lin}(\nu),C_{\rm lin}(\eta)\}.
 \label{eq:mhd-linear-constant}
\end{equation}
Let \(C_{\rm quad}^{\rm MHD}\) be a fixed dimension-dependent constant for
the sum of the four \(H^3\)-based quadratic estimates in
\eqref{eq:mhd-momentum} and \eqref{eq:mhd-induction}, and set
\begin{equation}
 c_{\nu,\eta}^{\rm MHD}
 :=\left[
  8C_{\rm quad}^{\rm MHD}C_{\rm sys}(\nu,\eta)^2
 \right]^{-2}.
 \label{eq:mhd-restart-constant}
\end{equation}
For data \((a,c)\in H^3_\sigma\times H^3_\sigma\), define the paired mild
map
\begin{align}
 \Phi_u(w,h)(t)
 &:={}e^{\nu t\Delta}a
 -\int_0^t e^{\nu(t-s)\Delta}\PP\left[
  (w\cdot\nabla)w-(h\cdot\nabla)h
 \right](s)\,ds,
 \notag\\
 \Phi_b(w,h)(t)
 &:={}e^{\eta t\Delta}c
 -\int_0^t e^{\eta(t-s)\Delta}\left[
  (w\cdot\nabla)h-(h\cdot\nabla)w
 \right](s)\,ds.
 \label{eq:mhd-mild-map}
\end{align}
Equip \(X_\delta\times X_\delta\), with \(X_\delta\) from
Appendix~\ref{app:local-theory}, with the sum norm, and write
\(A=\norm a_{H^3}+\norm c_{H^3}\).  The two linear estimates and the four
quadratic products give
\begin{align}
 \norm{(\Phi_u,\Phi_b)(w,h)}_{X_\delta\times X_\delta}
 &\le C_{\rm sys}A
 +C_{\rm sys}C_{\rm quad}^{\rm MHD}\delta^{1/2}
  \norm{(w,h)}_{X_\delta\times X_\delta}^2,
 \label{eq:mhd-map-bound}\\
 \norm{(\Phi_u,\Phi_b)(w,h)-(\Phi_u,\Phi_b)(v,k)}_{X_\delta\times X_\delta}
 &\le C_{\rm sys}C_{\rm quad}^{\rm MHD}\delta^{1/2}
 \left(\norm{(w,h)}_{X_\delta\times X_\delta}
       +\norm{(v,k)}_{X_\delta\times X_\delta}\right)
 \notag\\
 &\hspace{7em}\cdot
 \norm{(w-v,h-k)}_{X_\delta\times X_\delta}.
 \label{eq:mhd-map-difference}
\end{align}
Taking \(R=2C_{\rm sys}A\) and
\(4C_{\rm sys}C_{\rm quad}^{\rm MHD}\delta^{1/2}R\le1\) makes the
radius-\(R\) ball invariant and the map contractive.  Banach's theorem gives
\begin{equation}
 \delta(\nu,\eta,a,c)
 \ge\min\left\{1,
 c_{\nu,\eta}^{\rm MHD}
 \bigl(1+\norm a_{H^3}+\norm c_{H^3}\bigr)^{-2}
 \right\}.
 \label{eq:mhd-local-lifespan}
\end{equation}
At the common endpoint,
\(\norm{u(T)}_{H^3}+\norm{b(T)}_{H^3}
 \le\sqrt2\mathsf S_3^{\rm MHD}\).  Hence the same-history restart obeys
\begin{equation}
 \delta_T^{\rm MHD}
 \ge\min\left\{1,
 c_{\nu,\eta}^{\rm MHD}
 \bigl(1+\sqrt2\mathsf S_3^{\rm MHD}\bigr)^{-2}
 \right\}.
 \label{eq:mhd-restart}
\end{equation}
Suppose instead that the maximal smooth MHD time \(T_{\max}\) were finite.
Apply the estimates above on \([0,\tau]\), \(\tau<T_{\max}\).  Their uniform
right sides are independent of \(\tau\), and every integrated right side is
bounded by the same expression with \(T_{\max}\) in place of \(\tau\).
Monotone convergence gives every coupled rectangle on
\((0,T_{\max})\).  The endpoint moduli then produce one pair
\((u_*,b_*)\in H^\infty_\sigma\times H^\infty_\sigma\) and both normalized
pressure jets at \(T_{\max}\).  In particular,
\(\norm{u_*}_{H^3}+\norm{b_*}_{H^3}\le\sqrt2\mathsf S_3^{\rm MHD}\), so
\eqref{eq:mhd-restart} continues the pair beyond \(T_{\max}\).  Local
uniqueness identifies the new branch with the old one on their overlap.
Passing to the limit in \eqref{eq:mhd-velocity-recurrence} and
\eqref{eq:mhd-magnetic-recurrence} identifies every right row with its left
endpoint row.  The joined coupled history is smooth, contradicting finite
maximality.  This proves the theorem under the stated critical datum
condition.
\end{proof}

\section{Critical hyperdissipation in three dimensions}
\label{sec:hyperdissipative-application}

The axisymmetric argument changed the geometry seen by the stretching term
while leaving the Laplacian untouched.  We now keep unrestricted
three-dimensional transport and change the dissipation instead.  Consider
\begin{equation}
 \partial_tu+(u\cdot\nabla)u+\nabla p
 +\nu\Lambda^{2\alpha}u=0,
 \qquad \nabla\cdot u=0,
 \qquad \Lambda:=(-\Delta)^{1/2},
 \label{eq:hyper-system}
\end{equation}
on \(\R^3\), with \(\alpha\ge5/4\).  At the borderline value, the basic
energy dissipation carries exactly the spatial integrability needed to pay
the differentiated transport term.  Once that payment is made, the rest of
the proof can be read from the equation: the fractional operator changes the
width of each spatial rectangle, but it does not change compatibility,
pressure normalization, endpoint formation, or same-history restart.
The classical global regularity range begins at this exponent.
\bibref{TaoHyper}

\begin{theorem}[Quantitative hyperdissipative jet]
\label{thm:quantitative-hyperdissipative-jet}
Let \(\nu>0\), \(\alpha\ge5/4\), and
\(u_0\in\SSigma(\R^3)\).  Equation~\eqref{eq:hyper-system} has a unique
global smooth solution with normalized pressure
 \(p=R_iR_j(u_i u_j)\).  For every finite horizon \(T\), every finite
 temporal order, and every nonnegative real spatial order, its spatial rows, mixed
velocity and pressure rows, compatible fractional rectangles, endpoint
moduli, frequency tails, datum-depth ledger, and restart lifespan obey the
explicit estimates \eqref{eq:hyper-spatial-bound}--
\eqref{eq:hyper-restart-datum} below.
\end{theorem}

\begin{proof}
Put \(M_0:=\norm{u_0}_2\).  Pairing
\eqref{eq:hyper-system} with \(u\) gives
\begin{equation}
 \sup_{t\ge0}\norm{u(t)}_2\le M_0,
 \qquad
 \int_0^T\norm{\Lambda^\alpha u(t)}_2^2\,dt
 \le\frac{M_0^2}{2\nu}.
 \label{eq:hyper-base-energy}
\end{equation}
The transport term has vanished from this row, but it returns as soon as a
spatial derivative is applied.  The exponent \(5/4\) appears when we ask
whether the one integral in \eqref{eq:hyper-base-energy} is enough to absorb
that return at every order.

Choose the finite exponent
\begin{equation}
 q_\alpha:=
 \begin{cases}
  6/(3-2\alpha),&5/4\le\alpha<3/2,\\
  6,&\alpha\ge3/2,
 \end{cases}
 \qquad
 \vartheta_\alpha
 :=\frac3\alpha\left(\frac12-\frac1{q_\alpha}\right).
 \label{eq:hyper-q-choice}
\end{equation}
Then
\begin{equation}
 \frac3{q_\alpha}\le\alpha-1,
 \qquad 0<\vartheta_\alpha\le1.
 \label{eq:hyper-q-properties}
\end{equation}
The Gagliardo--Nirenberg inequality and
\eqref{eq:hyper-base-energy} first give
\begin{equation}
 \norm{u(t)}_{q_\alpha}
 \le G_\alpha M_0^{1-\vartheta_\alpha}
       \norm{\Lambda^\alpha u(t)}_2^{\vartheta_\alpha}.
 \label{eq:hyper-low-frequency-interpolation}
\end{equation}
Squaring, applying H\"older in time with exponents
\(1/\vartheta_\alpha\) and \(1/(1-\vartheta_\alpha)\), and then using
\eqref{eq:hyper-base-energy} yield the datum-side action
\begin{align}
 \int_0^T\norm{u(t)}_{q_\alpha}^2\,dt
 &\le\mathcal A_\alpha(M_0,\nu,T),
 \label{eq:hyper-lq-action}\\
 \mathcal A_\alpha(M_0,\nu,T)
 &:=G_\alpha^2M_0^2(2\nu)^{-\vartheta_\alpha}
       T^{1-\vartheta_\alpha}.
 \label{eq:hyper-lq-action-definition}
\end{align}

Fix \(s\ge\alpha-1\) and set \(r=s+1-\alpha\).  Moving \(\alpha\)
derivatives from the transport term onto the tested velocity gives
\begin{equation}
 \left|\pair{\Lambda^s\PP(u\cdot\nabla u)}{\Lambda^su}\right|
 \le
 \norm{\Lambda^r(u\otimes u)}_2
 \norm{\Lambda^{s+\alpha}u}_2.
 \label{eq:hyper-derivative-transfer}
\end{equation}
If \(p_\alpha=2q_\alpha/(q_\alpha-2)\), fractional Leibniz and homogeneous
Sobolev embedding give
\begin{align}
 \norm{\Lambda^r(u\otimes u)}_2
 &\le C_{\alpha,s}\norm u_{q_\alpha}
             \norm{\Lambda^ru}_{p_\alpha}
 \notag\\
 &\le C_{\alpha,s}\norm u_{q_\alpha}
             \norm{\Lambda^{r+3/q_\alpha}u}_2
 \le C_{\alpha,s}\norm u_{q_\alpha}\norm u_{H^s}.
 \label{eq:hyper-fractional-product}
\end{align}
The last inequality is exactly
\(r+3/q_\alpha\le s\), which is the first relation in
\eqref{eq:hyper-q-properties}.  Define
\begin{equation}
 X_s(t):=M_0^2+\norm{\Lambda^su(t)}_2^2,
 \qquad
 Y_s(t):=\norm{\Lambda^{s+\alpha}u(t)}_2.
 \label{eq:hyper-spatial-variables}
\end{equation}
Young's inequality now turns
\eqref{eq:hyper-derivative-transfer}--
\eqref{eq:hyper-fractional-product} into
\begin{equation}
 X_s'(t)+\nu Y_s(t)^2
 \le\frac{C_{\alpha,s}}{\nu}
       \norm{u(t)}_{q_\alpha}^2X_s(t).
 \label{eq:hyper-spatial-row}
\end{equation}
Set
\begin{equation}
 \Gamma_{\alpha,s}(T)
 :=\frac{C_{\alpha,s}}{\nu}
       \mathcal A_\alpha(M_0,\nu,T).
 \label{eq:hyper-spatial-exponent}
\end{equation}
Gronwall's inequality, followed once by the integrating-factor form of
\eqref{eq:hyper-spatial-row}, gives
\begin{align}
 \sup_{0\le t\le T}X_s(t)
 &\le X_s(0)e^{\Gamma_{\alpha,s}(T)},
 \label{eq:hyper-spatial-row-sup}\\
 \int_0^TY_s(t)^2\,dt
 &\le\frac{X_s(0)}{\nu}e^{\Gamma_{\alpha,s}(T)}.
 \label{eq:hyper-spatial-row-dissipation}
\end{align}
At \(\alpha=5/4\), one has \(q_\alpha=12\) and
\(\vartheta_\alpha=1\).  The derivative count in
\eqref{eq:hyper-fractional-product} becomes
\[
 s+1-\alpha+\frac3{q_\alpha}=s,
\]
and the exponent reduces to
\begin{equation}
 \Gamma_{5/4,s}(T)
 =\frac{C_sG_{5/4}^2M_0^2}{2\nu^2}.
 \label{eq:hyper-critical-exponent}
\end{equation}
No derivative and no time integrability remain unpaid at the critical
exponent.

For any requested real \(r\ge0\), put
\begin{equation}
 \bar s_r:=\max\{r,\alpha-1\},
 \qquad
 \kappa_{r,\alpha}
 :=\sup_{\rho\ge0}
 \frac{(1+\rho^2)^r}{1+\rho^{2\bar s_r}}<\infty,
 \qquad
 \mathsf S_r^\alpha(T)
 :=\kappa_{r,\alpha}^{1/2}X_{\bar s_r}(0)^{1/2}
    e^{\Gamma_{\alpha,\bar s_r}(T)/2}.
 \label{eq:hyper-spatial-envelope}
\end{equation}
Indeed,
\[
 \norm{u(t)}_{H^r}^2
 \le\kappa_{r,\alpha}\left(
  \norm{u(t)}_2^2+\norm{\Lambda^{\bar s_r}u(t)}_2^2
 \right)
 \le\kappa_{r,\alpha}X_{\bar s_r}(t).
\]
Thus
\begin{equation}
 \sup_{0\le t\le T}\norm{u(t)}_{H^r}
 \le\mathsf S_r^\alpha(T).
 \label{eq:hyper-spatial-bound}
\end{equation}
For the integrated row, let
\begin{equation}
 s_r^-:=\max\{\alpha-1,r-\alpha\}
 \label{eq:hyper-integrated-index}
\end{equation}
and define
\begin{equation}
 \bigl(\mathsf D_{0,r}^\alpha(T)\bigr)^2
 :=C_{\alpha,r}\left[
 TM_0^2+
 \frac{X_{s_r^-}(0)}{\nu}
 e^{\Gamma_{\alpha,s_r^-}(T)}
 \right].
 \label{eq:hyper-integrated-base}
\end{equation}
Since \(s_r^-+\alpha\ge r\), equations
\eqref{eq:hyper-base-energy} and
\eqref{eq:hyper-spatial-row-dissipation} give
\begin{equation}
 \norm u_{L^2(0,T;H^r)}
 \le\mathsf D_{0,r}^\alpha(T).
 \label{eq:hyper-integrated-base-bound}
\end{equation}
Choosing one \(s>\max\{5/2,\alpha-1\}\) in
\eqref{eq:hyper-spatial-row-sup} prevents a finite-time loss of the local
classical solution.  This already proves global smoothness.  We keep going
because the stronger conclusion is the complete quantitative jet.

\subsection{The fractional mixed and pressure jets}

Let \(V_n:=\partial_t^nu\) and \(Q_n:=\partial_t^np\).  Differentiating
\eqref{eq:hyper-system} and then applying \(\PP\) gives
\begin{equation}
 V_{n+1}
 =-\nu\Lambda^{2\alpha}V_n
 -\sum_{a=0}^n\binom na
   \PP(V_a\cdot\nabla V_{n-a}).
 \label{eq:hyper-velocity-recurrence}
\end{equation}
Taking divergence before projection fixes every pressure row:
\begin{equation}
 Q_n
 =R_iR_j\sum_{a=0}^n\binom na
   (V_a)_i(V_{n-a})_j.
 \label{eq:hyper-pressure-recurrence}
\end{equation}
For real \(r\ge0\), let \(\sigma_r:=\max\{2,r+1\}\), and enlarge the
fixed constant \(A_r\) so that
\begin{equation}
 \norm{\PP(f\cdot\nabla g)}_{H^r}
 +\norm{R_iR_j(f_i g_j)}_{H^r}
 \le A_r\norm f_{H^{\sigma_r}}\norm g_{H^{\sigma_r}}.
 \label{eq:hyper-row-product}
\end{equation}
Starting from
\begin{equation}
 \mathsf V_{0,r}^\alpha(T):=\mathsf S_r^\alpha(T),
 \label{eq:hyper-row-base}
\end{equation}
define
\begin{equation}
 \mathsf V_{n+1,r}^\alpha(T)
 :=\nu\mathsf V_{n,r+2\alpha}^\alpha(T)
 +A_r\sum_{a=0}^n\binom na
   \mathsf V_{a,\sigma_r}^\alpha(T)
   \mathsf V_{n-a,\sigma_r}^\alpha(T),
 \label{eq:hyper-row-recursion}
\end{equation}
and
\begin{equation}
 \mathsf Q_{n,r}^\alpha(T)
 :=A_r\sum_{a=0}^n\binom na
   \mathsf V_{a,\sigma_r}^\alpha(T)
   \mathsf V_{n-a,\sigma_r}^\alpha(T).
 \label{eq:hyper-pressure-envelope}
\end{equation}
Induction through \eqref{eq:hyper-velocity-recurrence} and
\eqref{eq:hyper-pressure-recurrence} proves
\begin{align}
 \sup_{0\le t\le T}\norm{V_n(t)}_{H^r}
 &\le\mathsf V_{n,r}^\alpha(T),
 \label{eq:hyper-row-bound}\\
 \sup_{0\le t\le T}\norm{Q_n(t)}_{H^r}
 &\le\mathsf Q_{n,r}^\alpha(T).
 \label{eq:hyper-pressure-bound}
\end{align}

The integrated recursion begins with
\eqref{eq:hyper-integrated-base}.  Define
\begin{equation}
 \mathsf D_{n+1,r}^\alpha(T)
 :=\nu\mathsf D_{n,r+2\alpha}^\alpha(T)
 +A_r\sum_{a=0}^n\binom na
   \mathsf V_{a,\sigma_r}^\alpha(T)
   \mathsf D_{n-a,\sigma_r}^\alpha(T).
 \label{eq:hyper-integrated-recursion}
\end{equation}
The same product estimate, with one factor in \(L_t^\infty\) and the other
in \(L_t^2\), gives
\begin{equation}
 \norm{V_n}_{L^2(0,T;H^r)}
 \le\mathsf D_{n,r}^\alpha(T).
 \label{eq:hyper-integrated-row-bound}
\end{equation}
The corresponding pressure row is
\begin{equation}
 \norm{Q_n}_{L^2(0,T;H^r)}
 \le\mathsf Q_{n,r}^{\alpha,(2)}(T)
 :=A_r\sum_{a=0}^n\binom na
   \mathsf V_{a,\sigma_r}^\alpha(T)
   \mathsf D_{n-a,\sigma_r}^\alpha(T).
 \label{eq:hyper-integrated-pressure}
\end{equation}

\subsection{Fractional rectangles and their common endpoint}

The dissipation selects the Gelfand triple
\begin{equation}
 H^{M+\alpha}_\sigma(\R^3)
 \hookrightarrow H^M_\sigma(\R^3)
 \hookrightarrow H^{M-\alpha}_\sigma(\R^3),
 \qquad M\ge0.
 \label{eq:hyper-gelfand-triple}
\end{equation}
For \(N\in\Nzero\) and real \(M\ge0\), define
\begin{equation}
 \mathcal N_{N,M}^\alpha(u;T)
 :=\sum_{n=0}^N\left[
  \norm{V_n}_{L^2(0,T;H^{M+\alpha})}^2
  +\norm{V_{n+1}}_{L^2(0,T;H^{M-\alpha})}^2
 \right].
 \label{eq:hyper-rectangle-norm}
\end{equation}
Write \(x_+:=\max\{x,0\}\).  Since \(H^0\hookrightarrow H^r\) for
\(r<0\), \eqref{eq:hyper-integrated-row-bound} gives
\begin{align}
 \mathcal N_{N,M}^\alpha(u;T)
 &\le\mathfrak R_{N,M}^\alpha(T),
 \label{eq:hyper-rectangle-bound}\\
 \mathfrak R_{N,M}^\alpha(T)
 &:=\sum_{n=0}^N\left[
  \bigl(\mathsf D_{n,M+\alpha}^\alpha(T)\bigr)^2
  +\bigl(\mathsf D_{n+1,(M-\alpha)_+}^\alpha(T)\bigr)^2
 \right].
 \label{eq:hyper-rectangle-envelope}
\end{align}

The pressure also fits the upper side of the rectangle.  The real-order
product theorem gives
\begin{equation}
 \norm{fg}_{H^M}
 \le C_{\alpha,M}
 \norm f_{H^{M+\alpha}}\norm g_{H^{M+\alpha}},
 \label{eq:hyper-rectangle-product}
\end{equation}
because \(M<2(M+\alpha)-3/2\).  Hence
\begin{equation}
 \norm{Q_n}_{L^1(0,T;H^M)}
 \le C_{\alpha,M}\sum_{a=0}^n\binom na
 \mathsf D_{a,M+\alpha}^\alpha(T)
 \mathsf D_{n-a,M+\alpha}^\alpha(T).
 \label{eq:hyper-rectangle-pressure}
\end{equation}
The differentiated transport term lies in
\(L^1(0,T;H^{M-\alpha})\) because
\begin{equation}
 \norm{\PP(V_a\cdot\nabla V_b)}_{H^{M-\alpha}}
 \le C_{\alpha,M}
 \norm{V_a}_{H^{M+\alpha}}
 \norm{V_b}_{H^{M+\alpha}}.
 \label{eq:hyper-lower-transport}
\end{equation}
The remaining lower-row terms obey
\begin{equation}
 \norm{\Lambda^{2\alpha}V_n}_{H^{M-\alpha}}
 \le\norm{V_n}_{H^{M+\alpha}},
 \qquad
 \norm{\nabla Q_n}_{H^{M-\alpha}}
 \le\norm{Q_n}_{H^M}.
 \label{eq:hyper-lower-linear-pressure}
\end{equation}
Thus every term in the differentiated momentum equation occupies the lower member of
\eqref{eq:hyper-gelfand-triple}.

On the Fourier side the pivot factorization is exact:
\begin{equation}
 \langle\xi\rangle^{2M}
 =\langle\xi\rangle^{M-\alpha}
  \langle\xi\rangle^{M+\alpha}.
 \label{eq:hyper-pivot-factorization}
\end{equation}
Apply this identity first to smooth temporal cutoffs and then pass by density.
For \(0\le s<t\le T\), the adjacent-row identity gives
\begin{equation}
 \left|
  \norm{V_n(t)}_{H^M}^2-\norm{V_n(s)}_{H^M}^2
 \right|
 \le2\norm{V_n}_{L^2(s,t;H^{M+\alpha})}
       \norm{V_{n+1}}_{L^2(s,t;H^{M-\alpha})}.
 \label{eq:hyper-adjacent-trace}
\end{equation}
It follows that \(V_n\in C([0,T];H^M)\) for each finite pair \((n,M)\).
Every rectangle is cut from the same velocity and the same Riesz-normalized
pressure, so deleting higher rows or lowering the Sobolev order leaves its
overlap unchanged.  Their intersection is one mixed jet.

The uniform next row sharpens the trace to the endpoint moduli
\begin{align}
 \norm{V_n(t)-V_n(T)}_{H^r}
 &\le(T-t)\mathsf V_{n+1,r}^\alpha(T),
 \label{eq:hyper-endpoint-modulus}\\
 \norm{Q_n(t)-Q_n(T)}_{H^r}
 &\le(T-t)\mathsf Q_{n+1,r}^\alpha(T).
 \label{eq:hyper-pressure-endpoint-modulus}
\end{align}
If \(s>k+3/2-3/b\), \(2\le b\le\infty\), and
\(1\le a\le\infty\), Sobolev embedding also gives
\begin{align}
 \norm{\nabla^kV_n}_{L^a(0,T;L^b)}
 &\le C_{s,k,b}T^{1/a}\mathsf V_{n,s}^\alpha(T),
 \label{eq:hyper-mixed-norm}\\
 \norm{\nabla^kQ_n}_{L^a(0,T;L^b)}
 &\le C_{s,k,b}T^{1/a}\mathsf Q_{n,s}^\alpha(T),
 \label{eq:hyper-pressure-mixed-norm}
\end{align}
where \(T^{1/\infty}=1\).  Define the sharp high-frequency multiplier by
\begin{equation}
 \widehat{P_{>K}f}(\xi)
 :=\mathbf1_{\{|\xi|>K\}}\widehat f(\xi).
 \label{eq:hyper-tail-projector}
\end{equation}
For \(K\ge1\) and \(s>k\),
\begin{align}
 \sup_{0\le t\le T}
 \norm{P_{>K}\nabla^kV_n(t)}_2
 &\le K^{k-s}\mathsf V_{n,s}^\alpha(T),
 \label{eq:hyper-frequency-tail}\\
 \sup_{0\le t\le T}
 \norm{P_{>K}\nabla^kQ_n(t)}_2
 &\le K^{k-s}\mathsf Q_{n,s}^\alpha(T).
 \label{eq:hyper-pressure-frequency-tail}
\end{align}

The datum order stays finite at every requested coordinate.  In real-order
notation the velocity row \((n,r)\) reads \(u_0\) no higher than
\begin{equation}
 d_{n,r}^{u,\alpha}
 :=\max\{\alpha-1,r+2\alpha n\},
 \label{eq:hyper-velocity-ledger}
\end{equation}
and the pressure row no higher than
\begin{equation}
 d_{n,r}^{p,\alpha}
 :=\max\{\alpha-1,\sigma_r+2\alpha n\}.
 \label{eq:hyper-pressure-ledger}
\end{equation}
The rectangle \((N,M)\) reads no higher than
\begin{equation}
 d_{N,M}^{\rm rect,\alpha}
 :=\max\left\{
  \alpha-1,
  M+\alpha+2\alpha N,
  (M-\alpha)_++2\alpha(N+1)
 \right\}.
 \label{eq:hyper-rectangle-ledger}
\end{equation}
Taking ceilings converts these values into a classical integer-derivative
ledger.

\subsection{The fractional restart}

Fix \(s_R>3/2\), \(s_R\ge\alpha-1\), and put
\(v=u(T)\), \(M_R=\norm v_{H^{s_R}}\).  The fractional heat semigroup
satisfies
\begin{equation}
 \norm{\Lambda e^{-\nu t\Lambda^{2\alpha}}F}_{H^{s_R}}
 \le c_\alpha(\nu t)^{-1/(2\alpha)}\norm F_{H^{s_R}},
 \qquad c_\alpha=(2\alpha e)^{-1/(2\alpha)}.
 \label{eq:hyper-semigroup}
\end{equation}
The operator in the mild equation consequently satisfies
\begin{equation}
 \norm{e^{-\nu t\Lambda^{2\alpha}}\PP\nabla\cdot F}_{H^{s_R}}
 \le C_{\PP}c_\alpha(\nu t)^{-1/(2\alpha)}
       \norm F_{H^{s_R}}.
 \label{eq:hyper-mild-operator}
\end{equation}
Since \(H^{s_R}\) is an algebra, the mild bilinear map on
\(C([0,\delta];H^{s_R})\) has norm at most
\begin{equation}
 K_{\alpha,s_R}\nu^{-1/(2\alpha)}
 \delta^{1-1/(2\alpha)},
 \qquad
 K_{\alpha,s_R}
 :=\frac{C_{s_R}c_\alpha}{1-1/(2\alpha)}.
 \label{eq:hyper-bilinear-constant}
\end{equation}
Here \(C_{s_R}\) includes \(C_{\PP}\) and the \(H^{s_R}\)-algebra
constant.
The ball of radius \(2M_R\) is invariant and contractive when
\begin{equation}
 8K_{\alpha,s_R}\nu^{-1/(2\alpha)}
 \delta^{1-1/(2\alpha)}M_R\le1.
 \label{eq:hyper-contraction-condition}
\end{equation}
For \(M_R>0\), the resulting lifespan is
\begin{equation}
 \delta_T^\alpha
 \ge c_{\alpha,s_R}
 \nu^{1/(2\alpha-1)}
 M_R^{-2\alpha/(2\alpha-1)},
 \qquad
 c_{\alpha,s_R}:=(8K_{\alpha,s_R})^{-2\alpha/(2\alpha-1)}.
 \label{eq:hyper-restart}
\end{equation}
Using \(M_R\le\mathsf S_{s_R}^\alpha(T)\) makes the restart entirely
datum-side:
\begin{equation}
 \delta_T^\alpha
 \ge c_{\alpha,s_R}
 \nu^{1/(2\alpha-1)}
 \bigl(\mathsf S_{s_R}^\alpha(T)\bigr)^{-2\alpha/(2\alpha-1)}.
 \label{eq:hyper-restart-datum}
\end{equation}
For zero endpoint data the zero solution is global.  Otherwise the local
solution from \(v\) agrees with the old history wherever both exist.  To make
the terminal passage explicit, suppose a maximal time \(T_{\max}<\infty\)
existed and apply every estimate above on \([0,\tau]\),
\(\tau<T_{\max}\).  Each right side is bounded by the same nondecreasing
envelope evaluated at \(T_{\max}\).  Monotone convergence supplies every
fractional rectangle on \((0,T_{\max})\), while
\eqref{eq:hyper-pivot-factorization} and
\eqref{eq:hyper-adjacent-trace} give one common \(H^\infty\) endpoint.
The uniform next-row bounds identify the complete left jet there.  Inserting
that endpoint into \eqref{eq:hyper-restart-datum} gives a positive interval
crossing \(T_{\max}\), and uniqueness glues it to the same solution.  Finite
maximality is impossible.  At \(\alpha=5/4\), the lower bound has the
concrete form
\[
 \delta_T^{5/4}
 \ge c_{5/4,s_R}\nu^{2/3}
 \bigl(\mathsf S_{s_R}^{5/4}(T)\bigr)^{-5/3}.
\]
This completes the hyperdissipative application.
\end{proof}

\section{Elliptic regularization of the time derivative}
\label{sec:voigt-application}

The last change acts in a different place.  Hyperdissipation strengthens the
operator applied to \(u\); the Navier--Stokes--Voigt equation places an
elliptic operator on \(u_t\).  Work on the three-dimensional torus
\(\mathbb T^3=(\R/2\pi\mathbb Z)^3\) with mean-zero divergence-free fields,
and set
\begin{equation}
 A:=-\PP\Delta,
 \qquad G_\ell:=(I+\ell^2A)^{-1},
 \qquad B(v,w):=\PP\nabla\cdot(v\otimes w).
 \label{eq:voigt-operators}
\end{equation}
For \(\ell,\nu>0\), the projected equation is
\begin{equation}
 (I+\ell^2A)u_t+\nu Au+B(u,u)=0,
 \qquad u(0)=u_0.
 \label{eq:voigt-system}
\end{equation}
The inverse \(G_\ell\) cancels the two derivatives in the viscous term and
gains one derivative over transport.  Temporal differentiation now costs no
additional spatial order.  That change is visible not just in existence, but
in the exact datum ledger of the entire mixed jet.  Global well-posedness and
long-time dynamics for this model are classical.
\bibref{KalantarovTiti}

\begin{theorem}[Quantitative Navier--Stokes--Voigt jet]
\label{thm:quantitative-voigt-jet}
Fix \(\ell,\nu>0\).  Every mean-zero
\(u_0\in H^\infty_\sigma(\mathbb T^3)\) generates a unique global smooth
solution of \eqref{eq:voigt-system}.  For every \(T>0\), every
\(n,r,N,M\in\Nzero\), every \(m\in\mathbb N\) with \(m\ge2\), and every
\(K\ge1\), the spatial
producer, mixed
velocity rows, pressure rows, compatible rectangles, filtered critical
height, frequency tails, endpoint moduli, restart lifespan, and time
analyticity radius obey \eqref{eq:voigt-spatial-bound}--
\eqref{eq:voigt-analytic-radius} below.  The datum order of a velocity or
pressure row is independent of its temporal height.
\end{theorem}

\begin{proof}
Write
\begin{equation}
 |v|_s:=\norm{A^{s/2}v}_2\quad(s\in\mathbb R),
 \qquad
 E_q(t):=|u(t)|_{q-1}^2+\ell^2|u(t)|_q^2
 \quad(q\ge1).
 \label{eq:voigt-energy-level}
\end{equation}
Throughout this section the mean-zero space \(H^q_\sigma(\mathbb T^3)\) is
equipped with the \(A\)-scale norm \(|\cdot|_q\).  Poincar\'e's inequality
makes it equivalent to the usual inhomogeneous Sobolev norm, while this choice
keeps every displayed constant tied to the equation's own Fourier multiplier.
Pairing \eqref{eq:voigt-system} with \(u\) gives the exact modified energy
identity
\begin{equation}
 E_1(t)+2\nu\int_0^t|u(s)|_1^2\,ds=E_1(0).
 \label{eq:voigt-base-energy}
\end{equation}
Put \(K_1:=E_1(0)\).  At the next order, pairing with \(Au\) gives
\[
 \frac12E_2'(t)+\nu|u|_2^2=-\pair{B(u,u)}{Au}.
\]
The three-dimensional interpolation estimate
\begin{equation}
 |\pair{B(u,u)}{Au}|
 \le C\norm u_6\norm{\nabla u}_3|u|_2
 \le C|u|_1^{3/2}|u|_2^{3/2}
 \label{eq:voigt-second-row-product}
\end{equation}
and Young's inequality yield
\begin{equation}
 E_2'(t)+\nu|u|_2^2
 \le C_2\nu^{-3}\ell^{-6}K_1^3.
 \label{eq:voigt-second-row}
\end{equation}
Define
\begin{equation}
 K_2(T):=E_2(0)+C_2\nu^{-3}\ell^{-6}K_1^3T.
 \label{eq:voigt-K2}
\end{equation}

For every integer \(q\ge3\), pair the equation with \(A^{q-1}u\).  The
periodic Sobolev algebra estimate gives
\begin{equation}
 |B(u,u)|_{q-2}\le C_q|u|_{q-1}^2.
 \label{eq:voigt-high-product}
\end{equation}
Hence
\begin{align}
 |\pair{B(u,u)}{A^{q-1}u}|
 &\le C_q|u|_{q-1}^2|u|_q
 \notag\\
 &\le\frac\nu2|u|_q^2+C_q\nu^{-1}|u|_{q-1}^4.
 \label{eq:voigt-high-row-estimate}
\end{align}
Starting from \eqref{eq:voigt-K2}, define recursively
\begin{equation}
 K_q(T):=E_q(0)+C_q\nu^{-1}\ell^{-4}
 \int_0^T K_{q-1}(t)^2\,dt
 \qquad(q\ge3).
 \label{eq:voigt-Kq}
\end{equation}
Induction gives
\begin{align}
 \sup_{0\le t\le T}E_q(t)&\le K_q(T),
 \label{eq:voigt-energy-envelope}\\
 \nu\int_0^T|u(t)|_q^2\,dt&\le K_q(T)
 \qquad(q\ge2),
 \label{eq:voigt-dissipation-envelope}
\end{align}
while \eqref{eq:voigt-base-energy} gives
\(\int_0^T|u|_1^2\le K_1/(2\nu)\).  In particular,
\begin{equation}
 \sup_{0\le t\le T}|u(t)|_q
 \le\ell^{-1}K_q(T)^{1/2}.
 \label{eq:voigt-spatial-bound}
\end{equation}
Every \(K_q\) is a computable nondecreasing polynomial in \(T\), with
\begin{equation}
 \deg K_q\le 2^{q-1}-1\qquad(q\ge2).
 \label{eq:voigt-polynomial-degree}
\end{equation}
Its datum dependence enters only through
\(E_j(0)=|u_0|_{j-1}^2+\ell^2|u_0|_j^2\), \(j\le q\).

The base phase space already explains globality.  The Fourier multipliers
satisfy
\begin{equation}
 \norm{G_\ell Av}_{H^1}\le\ell^{-2}\norm v_{H^1},
 \qquad
 \norm{G_\ell B(v,w)}_{H^1}
 \le C\ell^{-2}\norm v_{H^1}\norm w_{H^1}.
 \label{eq:voigt-H1-vector-field}
\end{equation}
Thus \eqref{eq:voigt-system} is an ordinary differential equation on
\(H^1_\sigma\).  Its local solution can end only if the \(H^1\) norm becomes
unbounded, and \eqref{eq:voigt-base-energy} prevents that.  The higher
polynomials then propagate every smooth datum order globally.

\subsection{A mixed jet with no temporal derivative loss}

Let \(V_n:=\partial_t^nu\).  Applying \(G_\ell\) after differentiating
\eqref{eq:voigt-system} gives
\begin{equation}
 V_{n+1}
 =-\nu G_\ell AV_n
 -G_\ell\sum_{a=0}^n\binom na B(V_a,V_{n-a}).
 \label{eq:voigt-row-recurrence}
\end{equation}
For \(r\in\Nzero\), put \(\sigma_r:=\max\{2,r-1\}\).  The same multipliers and
Sobolev multiplication give
\begin{align}
 |G_\ell Av|_r&\le\ell^{-2}|v|_r,
 \label{eq:voigt-linear-multiplier}\\
 |G_\ell B(v,w)|_r
 &\le C_r\ell^{-2}|v|_{\sigma_r}|w|_{\sigma_r}.
 \label{eq:voigt-bilinear-multiplier}
\end{align}
Define the spatial base values
\begin{equation}
 S_0^{\rm V}(T):=K_1^{1/2},
 \qquad
 S_r^{\rm V}(T):=\ell^{-1}K_r(T)^{1/2}\quad(r\ge1),
 \label{eq:voigt-spatial-base}
\end{equation}
and set \(U_{0,r}^{\rm V}:=S_r^{\rm V}\).  The row recursion is
\begin{equation}
 U_{n+1,r}^{\rm V}(T)
 :=\ell^{-2}\left[
  \nu U_{n,r}^{\rm V}(T)
  +C_r\sum_{a=0}^n\binom na
    U_{a,\sigma_r}^{\rm V}(T)
    U_{n-a,\sigma_r}^{\rm V}(T)
 \right].
 \label{eq:voigt-uniform-recursion}
\end{equation}
It follows that
\begin{equation}
 \sup_{0\le t\le T}|V_n(t)|_r
 \le U_{n,r}^{\rm V}(T).
 \label{eq:voigt-uniform-row}
\end{equation}
Unlike the Laplacian and fractional recursions, the right side of
\eqref{eq:voigt-uniform-recursion} never asks for a higher spatial order of
\(V_n\).  Consequently,
\begin{equation}
 U_{n,r}^{\rm V}(T)
 \text{ reads }u_0\text{ only through }H^{\max\{2,r\}},
 \label{eq:voigt-row-ledger}
\end{equation}
independently of \(n\).

The integrated rows are just as explicit.  Put
\begin{align}
 D_{0,0}^{\rm V}(T)&:=T^{1/2}S_0^{\rm V}(T),
 \notag\\
 D_{0,1}^{\rm V}(T)&:=\left(\frac{K_1}{2\nu}\right)^{1/2},
 \notag\\
 D_{0,r}^{\rm V}(T)&:=\left(\frac{K_r(T)}{\nu}\right)^{1/2}
 \quad(r\ge2).
 \label{eq:voigt-integrated-base}
\end{align}
Define
\begin{equation}
 D_{n+1,r}^{\rm V}(T)
 :=\ell^{-2}\left[
  \nu D_{n,r}^{\rm V}(T)
  +C_r\sum_{a=0}^n\binom na
    U_{a,\sigma_r}^{\rm V}(T)
    D_{n-a,\sigma_r}^{\rm V}(T)
 \right].
 \label{eq:voigt-integrated-recursion}
\end{equation}
Taking one factor in \(L_t^\infty\) and the other in \(L_t^2\) gives
\begin{equation}
 \norm{V_n}_{L^2(0,T;H^r)}
 \le 2^{r_+/2}D_{n,r}^{\rm V}(T),
 \qquad r_+:=\max\{r,0\}.
 \label{eq:voigt-integrated-row}
\end{equation}

For the lower side of the \(M=0\) rectangle, set
\begin{equation}
 \sigma_{-1}:=2,
 \quad U_{0,-1}^{\rm V}:=S_0^{\rm V},
 \quad D_{0,-1}^{\rm V}:=T^{1/2}S_0^{\rm V},
 \label{eq:voigt-negative-base}
\end{equation}
and extend \eqref{eq:voigt-uniform-recursion} and
\eqref{eq:voigt-integrated-recursion} to \(r=-1\) using
\begin{equation}
 |G_\ell Av|_{-1}\le\ell^{-2}|v|_{-1},
 \qquad
 |G_\ell B(v,w)|_{-1}
 \le C\ell^{-2}|v|_2|w|_2.
 \label{eq:voigt-negative-multipliers}
\end{equation}

\subsection{Pressure, rectangles, and the filtered critical height}

Restoring pressure gives
\begin{equation}
 (I-\ell^2\Delta)u_t-\nu\Delta u
 +(u\cdot\nabla)u+\nabla p=0.
 \label{eq:voigt-unprojected}
\end{equation}
The filtered time derivative remains divergence free, so
\begin{equation}
 Q_n:=\partial_t^np
 =R_iR_j\sum_{a=0}^n\binom na(V_a)_i(V_{n-a})_j.
 \label{eq:voigt-pressure-recurrence}
\end{equation}
With \(\tau_r:=\max\{2,r\}\), define
\begin{align}
 P_{n,r}^{\rm V}(T)
 &:=C_r\sum_{a=0}^n\binom na
  U_{a,\tau_r}^{\rm V}(T)U_{n-a,\tau_r}^{\rm V}(T),
 \label{eq:voigt-pressure-envelope}\\
 \Pi_{n,r}^{\rm V}(T)
 &:=C_r\sum_{a=0}^n\binom na
  U_{a,\tau_r}^{\rm V}(T)D_{n-a,\tau_r}^{\rm V}(T).
 \label{eq:voigt-integrated-pressure-envelope}
\end{align}
Then
\begin{align}
 \sup_{0\le t\le T}\norm{Q_n(t)}_{H^r}
 &\le P_{n,r}^{\rm V}(T),
 \label{eq:voigt-pressure-bound}\\
 \norm{Q_n}_{L^2(0,T;H^r)}
 &\le\Pi_{n,r}^{\rm V}(T).
 \label{eq:voigt-integrated-pressure-bound}
\end{align}
The pressure row reads no datum order above \(\max\{2,r\}\), again
independently of \(n\).

For \(N,M\in\Nzero\), let
\begin{equation}
 \mathcal N_{N,M}^{\rm V}(u;T)
 :=\sum_{n=0}^N\left[
  \norm{V_n}_{L^2(0,T;H^{M+1})}^2
  +\norm{V_{n+1}}_{L^2(0,T;H^{M-1})}^2
 \right].
 \label{eq:voigt-rectangle-norm}
\end{equation}
Equations \eqref{eq:voigt-integrated-row}--
\eqref{eq:voigt-negative-multipliers} give
\begin{align}
 \mathcal N_{N,M}^{\rm V}(u;T)
 &\le\mathfrak R_{N,M}^{\rm V}(T),
 \label{eq:voigt-rectangle-bound}\\
 \mathfrak R_{N,M}^{\rm V}(T)
 &:=\sum_{n=0}^N\left[
  2^{M+1}(D_{n,M+1}^{\rm V}(T))^2
  +2^{(M-1)_+}(D_{n+1,M-1}^{\rm V}(T))^2
 \right].
 \label{eq:voigt-rectangle-envelope}
\end{align}
The complete rectangle reads \(u_0\) only through
\begin{equation}
 H^{\max\{2,M+1\}},
 \label{eq:voigt-rectangle-ledger}
\end{equation}
independently of \(N\).  Compatibility is exact because every row and every
pressure is a restriction of one periodic solution.

The natural critical quantity for this equation includes its time filter:
\begin{equation}
 \mathcal H^{\rm V}(t)
 :=\frac12\left(|u(t)|_{1/2}^2+\ell^2|u(t)|_{3/2}^2\right).
 \label{eq:voigt-filtered-critical-height}
\end{equation}
Poincar\'e monotonicity and the first two energy envelopes give
\begin{equation}
 \sup_{0\le t\le T}\mathcal H^{\rm V}(t)
 \le\frac12\left(\ell^{-2}K_1+K_2(T)\right).
 \label{eq:voigt-filtered-critical-height-bound}
\end{equation}
Define the ordinary kinetic high-frequency tail by
\begin{equation}
 H_{>K}^{\rm V}(t)
 :=\frac12\sum_{|k|>K}|k|\,|\widehat u(k,t)|^2.
 \label{eq:voigt-frequency-tail-definition}
\end{equation}
Then
\begin{equation}
 \sup_{0\le t\le T}H_{>K}^{\rm V}(t)
 \le\frac12K^{1-2m}(S_m^{\rm V}(T))^2,
 \qquad m\in\mathbb N,\ m\ge2,\ K\ge1.
 \label{eq:voigt-frequency-tail}
\end{equation}

\subsection{Endpoint, restart, and time analyticity}

The uniform next row gives
\begin{align}
 |V_n(T)-V_n(t)|_r
 &\le(T-t)U_{n+1,r}^{\rm V}(T),
 \label{eq:voigt-endpoint-modulus}\\
 \norm{Q_n(T)-Q_n(t)}_{H^r}
 &\le(T-t)P_{n+1,r}^{\rm V}(T).
 \label{eq:voigt-pressure-endpoint-modulus}
\end{align}
Every finite rectangle reaches the same \(H^\infty\) endpoint and the same
normalized pressure jet.

On an \(H^1\) ball of radius \(R\), the vector field in
\eqref{eq:voigt-H1-vector-field} has Lipschitz constant at most
\(C\ell^{-2}(\nu+R)\).  Picard iteration, with the torus and Poincar\'e
constants absorbed into \(c_{\rm V},C_{\rm V}\), gives
\begin{equation}
 \delta(w)
 \ge c_{\rm V}\min\left\{1,
 \frac{\ell^2}{\nu+1+\norm w_{H^1}}
 \right\}.
 \label{eq:voigt-local-lifespan}
\end{equation}
Since \(\norm{u(T)}_{H^1}\le C_{\rm V}\ell^{-1}K_1^{1/2}\),
\begin{equation}
 \delta_T^{\rm V}
 \ge c_{\rm V}\min\left\{1,
 \frac{\ell^2}{\nu+1+C_{\rm V}\ell^{-1}K_1^{1/2}}
 \right\}.
 \label{eq:voigt-restart}
\end{equation}
The bound does not deteriorate with the horizon.  It crosses any alleged
finite maximal time, and uniqueness glues the restarted curve to the old
history.  Induction through \eqref{eq:voigt-row-recurrence} identifies the
entire right jet with the left endpoint jet.

The lossless recursion gives one further quantitative consequence.  At
\(r=2\), set
\begin{equation}
 R_T:=S_2^{\rm V}(T),
 \qquad
 L_T:=\ell^{-2}\bigl(\nu+C_2R_T\bigr).
 \label{eq:voigt-analytic-constants}
\end{equation}
Since
\begin{equation}
 \sum_{a=0}^n\binom na a!(n-a)!=(n+1)n!,
 \label{eq:voigt-binomial-factorial}
\end{equation}
induction in \eqref{eq:voigt-uniform-recursion} gives
\begin{equation}
 U_{n,2}^{\rm V}(T)\le n!R_TL_T^n.
 \label{eq:voigt-factorial-row}
\end{equation}
Taylor's formula with integral remainder and
\eqref{eq:voigt-factorial-row} give, whenever the segment from \(t\) to
\(t+h\) lies in \([0,T]\),
\begin{equation}
 \norm{
 u(t+h)-\sum_{j=0}^{N}\frac{h^j}{j!}V_j(t)
 }_{H^2}
 \le 2R_T(L_T|h|)^{N+1}.
 \label{eq:voigt-taylor-remainder}
\end{equation}
Hence the solution is real-time analytic as an \(H^2_\sigma\)-valued curve
with radius at least
\begin{equation}
 \rho_T^{\rm V}
 \ge\frac1{L_T}
 =\frac{\ell^2}{\nu+C_2\ell^{-1}K_2(T)^{1/2}}.
 \label{eq:voigt-analytic-radius}
\end{equation}
The analytic-ODE continuation supplies the same radius at interval endpoints.
The pressure recurrence has the same radius because its binomial sum is
bounded by \(C(n+1)!R_T^2L_T^n\).  This completes an all-data global
three-dimensional application in which the full temporal tower requires no
increase in datum depth.
\end{proof}

\section{What the applications have tested}
\label{sec:application-comparison}

The common construction has now survived four distinct changes to the
equation.  What remained fixed was not a particular estimate.  In each case a
datum-side producer first controlled every requested finite spatial row; the
actual differentiated equation then generated the temporal and pressure
rows; those rows agreed under restriction; their common endpoint satisfied
the limiting recurrence; and a quantitative local theorem restarted the same
history.  The table records the part that genuinely changed.

\begin{table}[H]
\centering
\small
\begin{tabular}{@{}
>{\raggedright\arraybackslash}p{0.20\textwidth}
>{\raggedright\arraybackslash}p{0.30\textwidth}
>{\raggedright\arraybackslash}p{0.19\textwidth}
>{\raggedright\arraybackslash}p{0.22\textwidth}@{}}
\toprule
Equation & Datum-side producer & Cost of one time row & Quantitative conclusion \\
\midrule
Axisymmetric no swirl
& Maximum principle for \(\omega^\theta/r\), followed by the stretching and logarithmic gradient bounds
& Two spatial derivatives
& All smooth data on \(\R^3\); complete jet and restart \\
Viscous--resistive MHD
& Small critical Besov norm and integrable coupled Lipschitz action
& Two spatial derivatives in both fields
& Small critical data on \(\R^3\); both field jets and both pressure conventions \\
Hyperdissipative flow
& Kinetic energy plus the sharp fractional product estimate
& \(2\alpha\) spatial derivatives
& All smooth data on \(\R^3\) for \(\alpha\ge5/4\); real-order rectangles \\
Navier--Stokes--Voigt
& Filtered energy hierarchy \(E_q\)
& No increase in spatial order
& All smooth periodic data; complete jet, restart, and time-analyticity radius \\
\bottomrule
\end{tabular}
\caption{The producer changes; finite compatibility, endpoint formation, and restart do not.}
\label{tab:application-comparison}
\end{table}

The applications reach farther than four global-existence statements.  They
show which parts of the proof are geometric, which depend on the order of
dissipation, which survive coupling, and which improve when the time
derivative is elliptically filtered.  In every case the output remains a
finite datum-resolved computation at each requested coordinate rather than an
unquantified appeal to smoothness.
